\chapter{Conclusion}\label{chap:conclusion}

In this thesis, we studied how risk-based portfolio optimization performs for an investor who holds assets denominated
in several currencies, and whether optimizing the asset and currency weights jointly improves on the standard
sequential currency overlay, in which the asset weights are chosen first and the currency positions afterwards. We
considered five risk measures, variance, semi-variance, mean absolute deviation, Expected Shortfall and second-order
Expected Shortfall, and investable universes of three sizes, built from the largest stocks of 14 economies
(\Cref{chap:data-and-implementation}). In \Cref{chap:backtest-analysis} we backtested the joint and overlay programs in
an unregularized, a regularized and a mean-risk formulation against the equally weighted portfolio and two asset-only
benchmarks, and in \Cref{chap:joint-vs-sequential} we compared how well the two programs minimize their target risk
measure in and out of sample. We find that the optimizers achieve a much lower risk than the equally weighted portfolio
in every universe and higher risk-adjusted returns in the medium and large universes, which a mean term extends to the
small universe for the deviation measures. The benefit of joint over sequential optimization is limited to two
settings, objectives with a mean term in the medium and large universes, and risk minimization in the small universe.

\section{Findings}

Out of sample, the risk-minimizing optimizers reduce volatility by $34\%$ to $64\%$ relative to the equally weighted
portfolio, and its maximum drawdown by up to $30$ percentage points (\Cref{sec:baseline-backtests}). Without a return
term in the objective and with a lower exposure to the common movement of the equity markets, they earn a smaller share
of the equity risk premium and lower excess returns. On a risk-adjusted basis, they still outperform the equally
weighted portfolio, with Sharpe ratios of $0.41$ to $0.56$ against $0.33$ in the medium universe and of $0.51$ to $0.76$
against $0.37$ in the large universe, but not consistently in the small universe, where the equally weighted portfolio
reaches $0.13$. This pattern is consistent with the low-volatility anomaly, the empirical finding that low-risk stocks
earn higher risk-adjusted returns than high-risk stocks, and with the evidence that minimum-risk portfolios can
outperform naive diversification out of sample. The optimizers tilt the portfolio toward low-risk stocks, and a larger
universe offers more of them, while in the small universe they give up more return than they remove risk.

Regularization raises the Sharpe ratios further, most in the large universe, where they increase by $0.17$ on average
(\Cref{sec:effects-of-regularization}). The penalties limit the extreme positions that the unregularized programs fit to
noise in the lookback window, in line with the estimation error of optimized portfolios growing with the number of
assets relative to the length of the window. On the currency side, the regularized Joint and Overlay optimizers perform
on par with the Full optimizer, which keeps the currency exposure fully hedged, so without a mean term, actively
managing the currencies adds little over full hedging. The Zero optimizer, which leaves the currency exposure unhedged,
outperforms them only in the small universe. In the medium and large universes, where currency risk makes up a larger
share of the portfolio's total risk, managing the currency exposure pays off.

Adding a mean term raises the Sharpe ratios in almost every case, by $0.16$ to $0.24$ on average for the deviation
measures and by much less for the tail measures (\Cref{sec:effects-of-adding-a-mean-term}). The deviation measures are
invariant to a shift in the portfolio return, so minimizing them ignores expected returns entirely, whereas Expected
Shortfall and second-order Expected Shortfall already reward a higher mean. Since the expected excess currency returns
are set to the interest rate differential, the mean term likely also tilts the currency positions toward
higher-yielding currencies, so that they earn part of the carry premium that pure risk minimization gives up. With the
mean term, the variance and semi-variance optimizers earn higher excess returns than the equally weighted portfolio in
every universe, at $50\%$ to $70\%$ of its volatility.

Comparing all backtests in
\Cref{tab:baseline,tab:regularized-top10,tab:regularized-top50,tab:regularized-top100,tab:mean-risk}, the mean-variance
Joint optimizer, that is, the mean-risk Joint optimizer with variance as the risk measure, performs best out of sample.
Its Sharpe ratio is the highest in the medium universe, at $0.74$ against $0.33$ for the equally weighted portfolio, and
in the large universe, at $0.91$ against $0.37$, where the regularized Joint optimizer under second-order Expected
Shortfall reaches the same value. In the small universe, its Sharpe ratio of $0.37$ against $0.13$ is exceeded only by
the mean-risk semi-variance optimizers at $0.38$. Its Sortino ratio confirms the ranking in the medium and large
universes, although several of the risk-minimizing optimizers have smaller drawdowns. The mean-variance Joint optimizer
combines two sources of improvement. Variance is a deviation measure, which gains the most from the mean term, and the
joint program can weigh the expected returns of assets and currencies against each other, which the overlay program
cannot since it fixes the asset weights first.

For the main question of the thesis, the comparison between the two approaches depends on whether expected returns
enter the objective. When the optimizers only minimize risk, the Joint and Overlay optimizers perform very similarly,
with or without regularization, and the most consistent difference is the lower volatility of the Joint optimizers,
although the difference shrinks as the universe grows. With a mean term, the Joint optimizers achieve higher Sharpe and
Sortino ratios than the Overlay for every risk measure in the medium and large universes, by about $0.07$ in Sharpe
ratio on average.

Since the lower volatility is the most consistent difference between the two approaches, \Cref{chap:joint-vs-sequential}
compares them on the risk measure they minimize. By construction, the joint program achieves an in-sample risk no higher
than that of the overlay program, since the solution of the overlay program is feasible for the joint program. Out of
sample, this advantage falls with universe size for all five risk measures. It is present only in the small universe,
where the currency weights are a larger share of the decision variables, and drops to zero in the larger universes,
where the asset step of the overlay program has enough weights to fit the objective on its own. Since the dependence
between asset and currency returns barely changes across universe sizes (\Cref{sec:data}), this decline is driven by the
number of assets relative to the number of currencies and not by a change in the dependence structure.

For the deviation measures, most of the small in-sample advantage is realized out of sample, while the larger in-sample
advantage of the tail measures is not. We trace it to two mechanisms that do not generalize, a tail-fitting mechanism,
in which the joint program aligns the currency weights with the specific worst days of the lookback window, and an
implicit risk-return tradeoff, through which the tail measures reward a high realized mean (\Cref{sec:mechanism-level}).
These results agree with the finding of \cite{RegularizedMultiCurrencyExpectedShortfallPortfolios} that the benefit of
joint optimization is concentrated in small universes, and qualify the finding of
\cite{SparseAndStableInternationalPortfolioOptimization} that joint optimization outperforms the overlay, since in our
backtests it does so in risk-adjusted terms only once a mean term enters the objective.

Overall, optimizing the portfolio pays off relative to naive diversification in the medium and large universes, which
offer more low-risk stocks to choose from, and more so with regularization and a mean term, which extends the advantage
to the small universe for the deviation measures. Joint optimization of the asset and currency weights is worth its
additional complexity in two situations, when the number of assets is small relative to the number of currencies, and
when expected returns enter the objective in the medium and large universes. When the goal is only to minimize risk in
larger universes, the sequential currency overlay performs just as well.

\section{Limitations and Future Research}

Our analysis covers one base currency, the US dollar, 14 economies, and one out-of-sample period from 2000 to 2023,
so the results reflect the dependence between asset and currency returns for a US investor over this period. For
an investor based in another currency, such as the Euro, the Swiss franc or the Japanese yen, the excess currency returns
and their dependence with the asset returns differ, and repeating the analysis for other base currencies would show how
much the results depend on this structure.

The backtests are also based on simplified trading assumptions, namely constant transaction costs, no market impact, and are conducted with quarterly
rebalancing, among others.
Since the unregularized and mean-risk Joint optimizers have a much higher currency turnover than the Overlay in most cases, higher trading
costs, more frequent rebalancing and market impact would likely reduce the advantage of the Joint optimizers
over the Overlay. A more realistic analysis could include a dynamic transaction cost model and 
multiple rebalancing frequencies, and test the sensitivity of the results to these assumptions.

Two further extensions follow from the results. First, the tail measures are estimated from the empirical distribution
of the lookback window, and are sensitive to estimation error. Estimators or formulations that
stabilize the tail, such as the extreme value approach proposed in \cite{McNeilFrey2000}, could improve the out-of-sample performance.
Second, the explanations given for the optimizers' performance, a lower market
exposure, a tilt towards low-risk stocks and currency carry, could be further investigated by
regressing their returns on market, low-volatility, carry and other factors.

